\section{Results}
% Present and synthesize the findings, tied to the research question.

\begin{figure}[t]
  \centering
  \includegraphics[width=0.495\linewidth]{attack_marital_education.png}\hfill
  \includegraphics[width=0.495\linewidth]{attack_marital_age.png}
  \caption{Sex reconstruction error under random tree selection, mean over 93--100
      valid seeds per condition and release size. Both panels compare None, A, B and
      A+B, with A = marital status. Left: B = education. Right: B = age.}
  \label{fig:attack}
\end{figure}

\begin{table}[t]
  \centering
  \caption{Sex error (\%), mean over seeds (94 at $m=200$). Gain: per-seed difference to
      None, averaged over all $m$. Holm-corrected $p$. None and A: education run.}
  \label{tab:attack}
  \small
  \setlength{\tabcolsep}{5pt}
  \begin{tabular}{lrrrrr}
    \toprule
    Known attributes (selection) & $m=1$ & $m=10$ & $m=200$ & Gain (points) & $p$ (Holm)
    \\
    \midrule
    None & 34.6 & 32.1 & 25.3 & -- & -- \\
    A: marital status & 34.2 & 31.3 & 24.4 & 0.5 & 0.036 \\
    B: education & 33.9 & 31.3 & 24.9 & 0.5 & 0.024 \\
    B: age & 34.5 & 31.8 & 25.1 & 0.1 & 0.42 \\
    A+B: marital status + education & 33.3 & 30.8 & 24.3 & 1.3 & $2\times10^{-6}$ \\
    A+B: marital status + age & 34.3 & 31.4 & 23.6 & 0.9 & 0.0015 \\
    A+B: marital status + education (sequential) & 33.7 & 33.6 & 30.4 & & \\
    \bottomrule
  \end{tabular}
\end{table}

\subsection{More trees, and a little background knowledge (SQ1)}
\label{sec:knowledge}
Figure~\ref{fig:attack} and Table~\ref{tab:attack} show that releasing more trees
has a larger effect than the tested background attributes. In the education run,
None error falls from 34.6\% to 25.3\%. Uniform and public-frequency controls have
50.4\% and 44.5\% error, illustrating why the choice of baseline matters.

A+B gives the largest average background gain, about one percentage point.
Education alone helps more than age alone, while age has no statistically significant
average gain. These comparisons use each run's own None control. Table~\ref{tab:attack}
reports seed-level corrected tests, but no per-$m$ test survives its separate Holm
correction. The small gains therefore need cautious interpretation.

Solver variability is comparable in size to these gains: identical None problems
sometimes return different errors across runs. Across shared cells, 15\% of None errors
differ, with a standard deviation of 2.0 percentage points for the paired differences.
Feasibility leaves the solver free to
choose among compatible tables. We also found no support for the proposed explanation
that gains track the share of trees splitting on sex and the known feature along the
same path. These checks limit how strongly we can explain the background effects.

\subsection{Which trees are released (SQ2)}
\label{sec:selection}
\begin{figure}[t]
  \centering
  \includegraphics[width=0.75\linewidth]{attack_selection_comparison.png}
  \caption{Sex reconstruction error with marital status and education known (A+B),
      comparing random tree selection with TreeFARMS enumeration order. Means use
      94--100 paired seeds per release size. Shaded bands show mean $\pm$ 1.96 SEM.}
  \label{fig:selection}
\end{figure}

With marital status and education known, random selection reaches 24.3\% error at
$m=200$, against 30.4\% for enumeration order (Figure~\ref{fig:selection}, paired
$p=8\times10^{-11}$). This larger difference shows that release policy matters under
the fixed attack. It does not by itself establish which structural property of the
selected trees causes the effect.
% AUDIT: the removed three-seed Jaccard check lacks a saved source artifact.

\begin{figure}[t]
  \centering
\includegraphics[width=0.37\linewidth]{confidence_reliability_m50.png}\hspace{0.05\linewidth}
  \includegraphics[width=0.37\linewidth]{confidence_public_data.png}
  \caption{Record-level confidence, 20 seeds. Left: stated confidence vs.\ share of
      correct claims at $m=50$ (marker area $\propto$ claims). Right: overconfidence of
      prior-guided sampling by public data and prior (95\% $t$-intervals).}
  \label{fig:confidence}
\end{figure}

\begin{table}[t]
  \centering
  \caption{Prior-guided sampling with 2,000 public records, means over 20 seeds.
      At $m=100$, 18 seeds returned candidates (1--20 tables each). Precision and AUROC
      are averaged per seed. Vouched: confidence $\geq0.9$.}
  \label{tab:confidence}
  \small
  \setlength{\tabcolsep}{4.5pt}
  \begin{tabular}{lrrrrr}
    \toprule
    & Accuracy (\%) & Overconfidence & Vouched (\%) & of which correct (\%) & AUROC \\
    \midrule
    Public statistics only & 58.5 & $+0.7$ & 10.3 & 92.0 & 0.74 \\
    $m=10$ & 69.2 & $+2.3$ & 30.9 & 95.3 & 0.79 \\
    $m=25$ & 73.9 & $+1.6$ & 39.3 & 95.4 & 0.81 \\
    $m=50$ & 77.8 & $+0.7$ & 45.0 & 97.6 & 0.83 \\
    $m=100$ & 78.9 & $+1.1$ & 47.1 & 96.8 & 0.82 \\
    \bottomrule
  \end{tabular}
\end{table}

\subsection{Record-level confidence and correctness (SQ3)}
\paragraph{Solver agreement alone is misleading.} Single-tree frequencies and
  repeated joint solves can assign high confidence to mostly incorrect claims. The
  solver can fill unconstrained cells similarly across solutions, so agreement is not
  evidence that those values are forced by the release.

\paragraph{Public guidance improves accuracy and mean confidence.} Prior-guided
  sampling is more accurate than the public-only control at every measured $m$
  (paired $p\leq2\times10^{-4}$). Its mean confidence exceeds accuracy by only
  0.7--2.3 points (Table~\ref{tab:confidence}). Figure~\ref{fig:confidence} left also
  compares scores with correctness. Small average gaps do not establish calibration
  at every confidence level.

\paragraph{Recovery is concentrated in high-confidence claims.} At $m=50$, claims
  with confidence $\geq0.9$ cover 45\% of the sample, with mean per-seed precision
  97.6\%. The public-only control covers about 10\%. This selective view shows
  exposure that a single overall accuracy hides, without identifying specific people.
  Even unanimous agreement is not certainty: at $m=50$, claims supported by all 20
  candidate tables were wrong in 2 of 673 cases.

\paragraph{A Dirichlet prior helps with a small public sample.} With 200 public
  records, plain-prior scores become overconfident (Figure~\ref{fig:confidence} right).
  At $m=10$, the Dirichlet variant reduces the mean gap from 8.6 to 1.6 points.
  The reduction is significant at all three tested sizes (paired absolute-gap
  $p\leq3\times10^{-5}$), with no measurable accuracy change. It selects fewer
  high-confidence claims but improves their precision. This suggests that uncertainty
  in public frequencies matters when little public data is available. With 2,000
  public records, the prior has no measurable effect on accuracy, ranking or the
  absolute mean confidence gap ($p\geq0.13$). Its benefit here is therefore tied
  to the small public sample, rather than a general improvement across settings.
